% =====================================================================
% 03_partie_a.tex — schéma, types, contraintes et données
% =====================================================================
\chapter{Partie A : schéma, contraintes et données}
\label{chap:partieA}

\section{Objectifs de la partie A}

La partie A établit le socle du projet. Elle crée l'utilisateur applicatif, les
types métier, les tables, les contraintes, les séquences et un jeu de données
représentatif, en suivant les conventions de SQL*Plus \cite{oracleSQLPlus}. Elle doit également garantir que l'ensemble peut être supprimé
puis recréé sans intervention manuelle. Cette exigence de rejeu est
particulièrement importante lors d'une soutenance : elle permet de repartir
d'un état propre si une démonstration a modifié les données.

\begin{objectifbox}[Résultat attendu]
À la fin de la partie A, la base contient 18 types métier, 11 tables logiques
(dont le stockage de la \emph{nested table} des signataires), 4 séquences et un
jeu de données couvrant les principaux cas de test. Aucun objet ne doit être
invalide.
\end{objectifbox}

\section{Les types métier}

Les types métier sont créés dans un ordre imposé par leurs dépendances. Un type
qui référence un autre type doit être créé après lui. Par exemple,
\code{client\_type} référence \code{adresse\_type}, \code{tel\_vry\_type} et
\code{conseiller\_type} ; il ne peut donc être déclaré avant eux. La
table~\ref{tab:types-metier} récapitule les dix-huit types de la partie A.

\begin{table}[htbp]
  \centering
  \small
  \caption{Types métier créés en partie A.}
  \label{tab:types-metier}
  \footnotesize
  \begin{tabularx}{\textwidth}{@{}L{5.4cm}L{3.1cm}X@{}}
    \toprule
    \entete{Type} & \entete{Nature} & \entete{Rôle} \\
    \midrule
    \code{adresse\_type} & Objet &
    \code{num\_rue}, \code{rue}, \code{ville}, \code{code\_postal}. \\
    \code{tel\_vry\_type} & \code{VARRAY(3)} &
    Téléphones d'un client. \\
    \code{conseiller\_type} & Objet &
    \code{num\_conseiller}, \code{nom}, \code{prenom}, \code{specialite}. \\
    \code{agence\_type} & Objet &
    \code{code}, \code{nom}, \code{ville}, \code{ref\_responsable}. \\
    \code{client\_type} & Objet &
    Identité, adresse, téléphones, conseiller, parrain. \\
    \code{signataire\_elt\allowbreak\_nt\_type} & Objet &
    \code{ref\_client REF client\_type}, \code{droit CHAR(1)}. \\
    \code{signataire\_nt\_type} & \emph{Nested table} &
    Collection de signataires. \\
    \code{compte\_type} & Objet, \code{NOT FINAL NOT INSTANTIABLE} &
    Super-classe des comptes. \\
    \code{cptCourant\_type} & Objet \code{UNDER compte\_type} &
    Découvert, \code{nb\_op\_cb}, signataires ; opérations via la table
    \code{operation}. \\
    \code{cptEpargne\_type} & Objet \code{UNDER compte\_type} &
    Taux d'intérêt, plafond de dépôt. \\
    \code{cptTerme\_type} & Objet \code{UNDER compte\_type} &
    Montant bloqué, durée, échéance, taux garanti. \\
    \code{operation\_type} & Objet &
    Opération sur un compte courant. \\
    \code{conseiller\_nt\_type} & \emph{Nested table} &
    Collection de conseillers pour la vue objet. \\
    \code{cptCourant\_nt\_type} & \emph{Nested table} &
    Collection de comptes courants pour la vue objet. \\
    \code{type\_stat\_elt\allowbreak\_nt\_type} & Objet &
    \code{type\_op}, \code{nb}, \code{montant\_total}. \\
    \code{type\_stat\_nt\_type} & \emph{Nested table} &
    Collection de statistiques. \\
    \code{mois\_elt\_vry\_type} & Objet &
    \code{mois}, \code{nb\_operations}, \code{types}. \\
    \code{mois\_vry\_type} & \code{VARRAY(12)} &
    Mois d'un relevé annuel. \\
    \code{agence\_vor\_type} & Objet (partie E) &
    Type de la vue objet \code{Agence\_VOR}. \\
    \bottomrule
  \end{tabularx}
\end{table}

\begin{remarquebox}[Pourquoi 26 lignes de type dans \code{USER\_OBJECTS} ?]
La partie A crée 18 types, et la partie E ajoute \code{agence\_vor\_type}, soit
19 types métier. Pourtant, la vue \code{USER\_OBJECTS} en affiche 26 : Oracle
crée des types internes \code{SYS\_YOID...} et conserve une version
supplémentaire de \code{CLIENT\_TYPE} après l'évolution
\code{ALTER TYPE ... ADD ATTRIBUTE}. Ce décalage n'est pas une anomalie ; il
provient du catalogue.
\end{remarquebox}

\section{Création de l'utilisateur et des privilèges}

Le script \code{00\_creation\_utilisateur.sql} doit être exécuté par un
administrateur, et non par \code{MASTER1}. Il crée l'utilisateur s'il n'existe
pas, puis accorde les privilèges nécessaires : connexion, création de tables,
de types, de vues, de procédures, de séquences et de triggers. Il accorde
également les droits sur \code{DBMS\_SCHEDULER} et un quota illimité sur
l'espace de travail.

\begin{lstlisting}[caption={Création idempotente de l'utilisateur \code{MASTER1}.},label={lst:creation-utilisateur}]
DECLARE
  v_nb NUMBER;
BEGIN
  SELECT COUNT(*) INTO v_nb FROM dba_users WHERE username = 'MASTER1';
  IF v_nb = 0 THEN
    EXECUTE IMMEDIATE 'CREATE USER MASTER1 IDENTIFIED BY system';
    DBMS_OUTPUT.PUT_LINE('Utilisateur MASTER1 cree.');
  ELSE
    DBMS_OUTPUT.PUT_LINE('Utilisateur MASTER1 deja present : aucune creation.');
  END IF;
END;
/

GRANT CREATE SESSION   TO MASTER1;
GRANT CREATE TABLE     TO MASTER1;
GRANT CREATE TYPE      TO MASTER1;
GRANT CREATE VIEW      TO MASTER1;
GRANT CREATE PROCEDURE TO MASTER1;
GRANT CREATE SEQUENCE  TO MASTER1;
GRANT CREATE TRIGGER   TO MASTER1;
GRANT CREATE JOB       TO MASTER1;
GRANT EXECUTE ON DBMS_SCHEDULER TO MASTER1;
GRANT UNLIMITED TABLESPACE TO MASTER1;
ALTER USER MASTER1 QUOTA UNLIMITED ON USERS;
\end{lstlisting}

\begin{attentionbox}[Sécurité]
Le mot de passe \code{system} est utilisé ici uniquement parce que
l'environnement est local et pédagogique. Dans un contexte réel, il faudrait un
mot de passe robuste, des profils de ressources, une séparation des rôles et
une journalisation des connexions. Cette limite est reprise dans les
perspectives.
\end{attentionbox}

\section{Les tables}

Le schéma comporte onze tables logiques, auxquelles s'ajoute la table physique
\code{signataire\_nt\_st} créée par Oracle pour stocker la \emph{nested table}
des signataires. Le tableau~\ref{tab:tables} présente leur nature et leur
contenu attendu.

\begin{table}[htbp]
  \centering
  \small
  \caption{Tables du schéma et contenu attendu.}
  \label{tab:tables}
  \begin{tabularx}{\textwidth}{@{}L{3.4cm}L{3.8cm}X@{}}
    \toprule
    \entete{Table} & \entete{Type} & \entete{Contenu} \\
    \midrule
    \code{conseiller} & Objet & 4 conseillers. \\
    \code{agence} & Objet & 3 agences. \\
    \code{affectation} & Relationnelle & Historique conseiller/agence. \\
    \code{client} & Objet & 6 clients. \\
    \code{cpt\_courant} & Objet &
    7 comptes courants au départ, 6 après les tests destructifs. \\
    \code{cpt\_epargne} & Objet &
    5 comptes d'épargne au départ, 4 après \code{fermeCompte('CE4')}. \\
    \code{cpt\_terme} & Objet &
    4 comptes à terme au départ, 3 après la requête 6.25. \\
    \code{operation} & Objet &
    20 opérations au départ, 15 après la suppression de CC1. \\
    \code{alerte} & Relationnelle &
    Alertes d'échéance et de solde. \\
    \code{"AUDIT"} & Relationnelle &
    Journal des modifications de taux et de découvert. \\
    \code{releves} & Relationnelle + \code{VARRAY} &
    Relevé annuel multi-niveaux. \\
    \code{signataire\_nt\_st} & Stockage de \emph{nested table} &
    Table physique créée par \code{NESTED TABLE ... STORE AS}. \\
    \bottomrule
  \end{tabularx}
\end{table}

\begin{attentionbox}[Le mot réservé \code{AUDIT}]
\code{AUDIT} est un mot réservé d'Oracle. Sans guillemets,
\code{CREATE TABLE AUDIT} échoue avec \code{ORA-00903}. La table est donc
toujours écrite \code{"AUDIT"}, à la création comme dans les insertions et les
requêtes. Cette contrainte de nommage est volontairement conservée pour
montrer la maîtrise des identifiants délimités.
\end{attentionbox}

\section{Types, tables et contraintes : extraits significatifs}

Le type \code{client\_type} illustre à lui seul plusieurs mécanismes : attributs
structurés, collection \code{VARRAY}, référence vers un conseiller et évolution
par \code{ALTER TYPE}. Le parrainage est ajouté après la création initiale, car
le sujet le demande en partie E tout en l'utilisant dès la partie A pour
insérer deux clients parrainés.

\begin{lstlisting}[caption={Création de \code{client\_type} et ajout du parrainage.},label={lst:client-type}]
BEGIN
  EXECUTE IMMEDIATE q'[
    CREATE TYPE client_type AS OBJECT (
      num_client     NUMBER,
      nom            VARCHAR2(30),
      prenom         VARCHAR2(30),
      date_naissance DATE,
      num_secu       VARCHAR2(20),
      adresse        adresse_type,
      telephones     tel_vry_type,
      email          VARCHAR2(60),
      date_entree    DATE,
      ref_conseiller REF conseiller_type
    )]';
END;
/

ALTER TYPE client_type ADD ATTRIBUTE ref_parrain REF client_type CASCADE;
\end{lstlisting}

Le recours à \code{EXECUTE IMMEDIATE} n'est pas un caprice de style. Sur
l'installation utilisée, un \code{CREATE TYPE} direct dans SQL*Plus ferme
silencieusement la session. Le bloc PL/SQL avec \code{q'[...]'} contourne ce
comportement et permet de vérifier les erreurs avec \code{SHOW ERRORS}.

La table \code{cpt\_courant} montre la contrainte différée utilisée par B5b.
Elle est bien \code{NOT NULL} sur le titulaire, mais la vérification est
repoussée au \code{COMMIT}. La \emph{nested table} des signataires est stockée
dans une table dédiée, sur laquelle des contraintes de domaine sont ajoutées.

\begin{lstlisting}[caption={Table \code{cpt\_courant} et stockage des signataires.},label={lst:cpt-courant}]
CREATE TABLE cpt_courant OF cptCourant_type (
  num_compte PRIMARY KEY,
  CONSTRAINT ck_cc_client    CHECK (ref_client IS NOT NULL)
                             DEFERRABLE INITIALLY DEFERRED,
  CONSTRAINT ck_cc_agence    CHECK (ref_agence IS NOT NULL),
  CONSTRAINT ck_cc_statut    CHECK (statut IN ('ACTIF','BLOQUE','DORMANT')),
  CONSTRAINT ck_cc_decouvert CHECK (decouvert >= 0),
  CONSTRAINT ck_cc_nbcb      CHECK (nb_op_cb >= 0)
)
NESTED TABLE signataires STORE AS signataire_nt_st;

ALTER TABLE signataire_nt_st
  ADD CONSTRAINT ck_sig_droit CHECK (droit IN ('X','R','D'));
ALTER TABLE signataire_nt_st
  ADD CONSTRAINT ck_sig_droit_nn CHECK (droit IS NOT NULL);
\end{lstlisting}

\section{Contraintes et valeurs par défaut}

Les contraintes ne se limitent pas aux clés primaires. Elles encodent des
règles métier : un compte doit avoir un titulaire, un statut appartient à un
domaine fermé, un taux d'épargne reste strictement inférieur à 3,5~\%, une
durée de compte à terme appartient à l'ensemble $\{6, 12, 24, 36\}$ mois, et un
découvert ne peut pas être négatif. Le tableau~\ref{tab:contraintes} en donne
une vue synthétique.

\begin{table}[htbp]
  \centering
  \small
  \caption{Contraintes et valeurs par défaut.}
  \label{tab:contraintes}
  \begin{tabularx}{\textwidth}{@{}L{5.2cm}L{3.4cm}X@{}}
    \toprule
    \entete{Contrainte} & \entete{Table} & \entete{Rôle} \\
    \midrule
    PK sur \code{num\_client}, \code{num\_conseiller}, \code{code}, etc. &
    Toutes & Identifiants. \\
    \code{UNIQUE(num\_secu)} & \code{client} &
    Un numéro de sécurité sociale unique. \\
    \code{CHECK(ref\_client IS NOT NULL)} &
    \code{cpt\_courant}, \code{cpt\_epargne}, \code{cpt\_terme} &
    Un compte doit avoir un titulaire. Sur \code{cpt\_courant} seulement, la
    contrainte est \code{DEFERRABLE INITIALLY DEFERRED} pour B5b. \\
    \code{CHECK(droit IN ('X',\allowbreak'R',\allowbreak'D'))} &
    \code{signataire\_nt\_st} &
    Domaine des droits. \\
    \code{CHECK(statut IN ('ACTIF',\allowbreak'BLOQUE',\allowbreak'DORMANT'))}
    & Comptes &
    Domaine du statut. \\
    \code{CHECK(type\_op IN ('DEPOT',\allowbreak'RETRAIT',\allowbreak
    'VIREMENT',\allowbreak'CB'))} &
    \code{operation} & Domaine du type d'opération. \\
    \code{CHECK(taux\_interet >= 0 AND taux\_interet < 3.5)} &
    \code{cpt\_epargne} & Limite du taux d'épargne. \\
    \code{CHECK(duree\_mois IN (6,12,24,36))} & \code{cpt\_terme} &
    Durées autorisées. \\
    \code{CHECK(decouvert >= 0)} & \code{cpt\_courant} &
    Découvert positif ou nul. \\
    \code{DEFAULT SYSDATE} & \code{operation.\allowbreak date\_op} &
    Date d'opération par défaut. \\
    \code{DEFAULT 'ACTIF'} & Statut des comptes &
    Statut par défaut. \\
    FK sur les \code{REF} & Comptes, opérations, agence &
    Intégrité référentielle. \\
    \bottomrule
  \end{tabularx}
\end{table}

\begin{clebox}[La contrainte différée de B5b]
La contrainte \code{ck\_cc\_client} est \code{DEFERRABLE INITIALLY DEFERRED}.
Elle interdit toujours un compte sans titulaire, mais la vérification a lieu au
\code{COMMIT} et non à l'insertion. C'est ce qui permet d'insérer \code{CC99}
sans titulaire, puis de le mettre à jour avant validation. Les contraintes
équivalentes de \code{cpt\_epargne} et \code{cpt\_terme} restent immédiates.
\end{clebox}

\section{Séquences}

Quatre séquences fournissent les identifiants des objets qui ne possèdent pas
de clé naturelle. Les opérations et les affectations commencent à 100, afin de
distinguer immédiatement les identifiants générés des petits identifiants de
test.

\begin{table}[htbp]
  \centering
  \small
  \caption{Séquences du schéma.}
  \label{tab:sequences}
  \begin{tabularx}{0.75\textwidth}{@{}L{4.2cm}X@{}}
    \toprule
    \entete{Séquence} & \entete{Utilisation} \\
    \midrule
    \code{seq\_operation} & Identifiants d'opération. \\
    \code{seq\_alerte} & Identifiants d'alerte. \\
    \code{seq\_audit} & Identifiants d'audit. \\
    \code{seq\_affectation} & Identifiants d'affectation. \\
    \bottomrule
  \end{tabularx}
\end{table}

\section{Jeu de données initial}

Le jeu de données est conçu pour couvrir les cas limites : clients avec zéro,
un, deux ou trois téléphones ; compte courant bloqué ; compte d'épargne proche
du taux maximal ; compte à terme dont l'échéance est proche, dépassée ou très
ancienne ; opérations de montants et de types variés. Le
tableau~\ref{tab:jeu-donnees} en donne la synthèse.

\begin{table}[htbp]
  \centering
  \small
  \caption{Jeu de données initial de la partie A.}
  \label{tab:jeu-donnees}
  \begin{tabularx}{\textwidth}{@{}L{3.3cm}C{1.6cm}X@{}}
    \toprule
    \entete{Élément} & \entete{Quantité} & \entete{Détail} \\
    \midrule
    Agences & 3 &
    A1 Béjaïa Centre, A2 Alger Grande Poste, A3 Oran Front de Mer. \\
    Conseillers & 4 &
    BENALI Karim, HAMDI Sara, TOUATI Yacine, KACI Nadia. \\
    Clients & 6 &
    AMRANI Ali, BENALI Samia, CHERIF Mohamed, DJEBBAR Lina, HADDAD Omar,
    ZERROUKI Nour. \\
    Comptes courants & 7 &
    CC1 à CC7, dont CC5 \code{BLOQUE}. \\
    Comptes d'épargne & 5 &
    CE1 à CE5, dont CE2 à 3,49~\% (proche de la limite de 3,5~\%). \\
    Comptes à terme & 4 &
    CT1 à CT4 ; CT4 a un montant nul et une échéance de plus d'un an pour la
    suppression 6.25. \\
    Opérations & 20 &
    Dépôts, retraits, virements et opérations CB. \\
    Signataires & 10 &
    Droits \code{X}, \code{R} et \code{D} sur au moins quatre comptes. \\
    \bottomrule
  \end{tabularx}
\end{table}

\begin{remarquebox}[Pourquoi un quatrième compte à terme ?]
\code{CT4} est un cas de test volontairement prévu pour la requête 6.25 :
\code{montant\_bloque = 0} et échéance de plus d'un an. Les autres comptes à
terme ont un montant bloqué strictement positif et ne doivent pas être
supprimés. Sans \code{CT4}, la requête de suppression ne démontrerait rien.
\end{remarquebox}

Le chargement des clients illustre la variété des \code{VARRAY} : le client 1
possède trois téléphones, le client 2 en possède deux avec un \code{NULL} au
milieu, le client 3 un seul téléphone non nul, et le client 5 un seul
téléphone. Cette variété est exploitée par la procédure
\code{change\_Portable} de la partie B.

\section{Rejeu et purge}

Le début de \code{partie\_A\_schema.sql} supprime les objets dans l'ordre
inverse des dépendances : vue, tables, séquences, puis types. Cet ordre est
indispensable ; supprimer un type encore référencé par une table produit une
erreur. Le script utilise \code{CASCADE CONSTRAINTS PURGE} pour les tables et
\code{FORCE} pour les types.

\begin{lstlisting}[caption={Nettoyage initial du schéma.},label={lst:purge-schema},basicstyle=\ttfamily\scriptsize]
DROP VIEW Agence_VOR;
DROP TABLE releves CASCADE CONSTRAINTS PURGE;
DROP TABLE "AUDIT" CASCADE CONSTRAINTS PURGE;
DROP TABLE alerte CASCADE CONSTRAINTS PURGE;
DROP TABLE operation CASCADE CONSTRAINTS PURGE;
DROP TABLE cpt_terme CASCADE CONSTRAINTS PURGE;
DROP TABLE cpt_epargne CASCADE CONSTRAINTS PURGE;
DROP TABLE cpt_courant CASCADE CONSTRAINTS PURGE;
DROP TABLE client CASCADE CONSTRAINTS PURGE;
DROP TABLE affectation CASCADE CONSTRAINTS PURGE;
DROP TABLE agence CASCADE CONSTRAINTS PURGE;
DROP TABLE conseiller CASCADE CONSTRAINTS PURGE;
DROP SEQUENCE seq_operation;
DROP SEQUENCE seq_alerte;
DROP SEQUENCE seq_audit;
DROP SEQUENCE seq_affectation;
DROP TYPE agence_vor_type FORCE;
DROP TYPE mois_vry_type FORCE;
DROP TYPE mois_elt_vry_type FORCE;
DROP TYPE type_stat_nt_type FORCE;
DROP TYPE type_stat_elt_nt_type FORCE;
DROP TYPE cptCourant_nt_type FORCE;
DROP TYPE conseiller_nt_type FORCE;
DROP TYPE operation_type FORCE;
DROP TYPE cptTerme_type FORCE;
DROP TYPE cptEpargne_type FORCE;
DROP TYPE cptCourant_type FORCE;
DROP TYPE compte_type FORCE;
DROP TYPE signataire_nt_type FORCE;
DROP TYPE signataire_elt_nt_type FORCE;
DROP TYPE client_type FORCE;
DROP TYPE agence_type FORCE;
DROP TYPE conseiller_type FORCE;
DROP TYPE tel_vry_type FORCE;
DROP TYPE adresse_type FORCE;
\end{lstlisting}

Deux scripts de maintenance complètent ce dispositif.
\code{purge\_donnees.sql} vide les données avec des \code{DELETE} contrôlés,
sans supprimer la structure. Le script unique
\code{projet\_bdda\_banque.sql} est une concaténation physique des parties A à
E, et non un script d'inclusion \code{@}. Cette choice évite les problèmes de
chemin et permet de déposer un seul fichier exécutable.

\begin{attentionbox}[Le script unique n'est pas un script de production]
Le script supprime et recrée tout. Il est conçu pour un TP et une démonstration,
pas pour une base en exploitation. Sur une base réelle, on utiliserait des
migrations versionnées, une sauvegarde préalable et des scripts de retour
arrière. La purge séparée permet de vider les données sans toucher au schéma.
\end{attentionbox}

\section{Points de contrôle et pièges}

La partie A doit être validée avant de passer à la suite. Les contrôles
portent sur l'absence d'objets invalides, la présence des contraintes, le
nombre de lignes et la possibilité de rejouer le script. Les captures
n°~01 à n°~11 illustrent ces étapes : connexion, création de l'utilisateur,
types, tables, contraintes, séquences, données et contrôle final.

\begin{table}[htbp]
  \centering
  \small
  \caption{Pièges rencontrés dans la partie A et solutions.}
  \label{tab:pieges-partieA}
  \begin{tabularx}{\textwidth}{@{}L{4.7cm}L{5.1cm}X@{}}
    \toprule
    \entete{Piège} & \entete{Symptôme} & \entete{Solution} \\
    \midrule
    \code{CREATE TYPE} direct dans SQL*Plus &
    La session se ferme silencieusement. &
    Utiliser \code{BEGIN EXECUTE IMMEDIATE q'[...]'; END; /}. \\
    Ligne terminée par un tiret &
    SQL*Plus la prend pour une continuation. &
    Éviter les lignes qui finissent par un tiret. \\
    \code{AUDIT} sans guillemets &
    \code{ORA-00903 : nom de table non valide}. &
    Écrire \code{"AUDIT"}. \\
    \code{ALTER TABLE} pour ajouter un attribut à une table objet &
    \code{ORA-22856}. &
    Utiliser \code{ALTER TYPE ... ADD ATTRIBUTE ... CASCADE}. \\
    Contrainte de même nom sur deux tables &
    \code{ORA-02264}. &
    Les noms de contraintes sont uniques dans le schéma ; préfixer. \\
    \code{REF} basé sur la PK dans une \emph{nested table} &
    \code{ORA-22979}. &
    Garder l'OID système pour \code{client}. \\
    \code{CARDINALITY} sur un \code{VARRAY} &
    \code{ORA-00932}. &
    Utiliser \code{TABLE(...)} et \code{COUNT(*)}. \\
    \bottomrule
  \end{tabularx}
\end{table}

\section{Conclusion de la partie A}

La partie A fournit un schéma complet, contraint et rejouable. Elle ne se
contente pas de créer des tables : elle encode les règles métier, prépare les
collections de la partie E et met en place les mécanismes de rejeu. Les
contraintes différées, les OID et les types récursifs sont déjà présents, ce
qui permet aux parties suivantes de se concentrer sur les manipulations et les
traitements.
